\onecolumn

\section{Admission test and ledger invariants}

{\small
\begin{verbatim}
State, per pool p: reserved set R_p with bound K_p; capacity vector Cap_p
(compute, device memory, spill-tier memory and restore bandwidth); forecast
window W; committed(p, t) = sum over r in R_p of r's remaining claim
profile over [t, t+W]; per-class budget quotas Q_c; per-tenant envelopes
E_g.

ADMIT(f) for flow f with descriptor (class c, tenant g, budget b, footprint
m, profile est):
  1. |R_p| < K_p
  2. committed(p, t) + claim(est) <= ceiling(p) for all t in window
  3. exists device set D: pinnable(m, D) under embodiment (strict: resident
     throughout; tiered: resident-active plus reserved spill and reserved
     restore within bound L_restore)
  4. spent(c) + b <= Q_c and within E_g
  5. trust(f) >= eligibility threshold for reservation
  All pass: grant r(f), append to ledger, R_p += {f}.
  Any fail: route to class c best-effort, or decline with optional quote.

Events: RENEW is ADMIT at current state with no carried privilege. EXPIRE
and EXHAUST release r(f) with the class's disposition. HALT(authority)
releases r(f) immediately; it is the only release not initiated by the
flow's own lifecycle.

Invariants (any interleaving of admit, renew, expire, exhaust, halt):
  I1 Compute: committed(p, t) <= ceiling(p) <= Cap_p at all t.
  I2 Memory: pinned bytes + reserved spill/restore never exceed their
     tiers; a tiered-pin flow always has a reserved restore path.
  I3 Budget conservation: metered decode tokens of f never exceed b; quota
     accounting is exact under renegotiation.
  I4 Guarantee monotonicity: once granted, r(f) is removed only by expire,
     exhaust, or halt; no capacity condition appears in any release path.
\end{verbatim}
}
